\section{本論文の目的}
\label{sec:purpose}

本論文では，嶋田義皓『量子コンピューティング
基本アルゴリズムから量子機械学習まで』\cite{shimada2020}を主資料とし，
量子計算の基礎からHHLアルゴリズムまでを扱う．
主資料の第2章から量子状態，測定，量子ゲートと複合系を，
第3章からHHLの中で使う個々の処理を取り上げ，
第4章4.4.2節の量子線形方程式へつなげる．
主資料で省略された途中計算は，どの線形代数の結果を使ったかを書いて補う．
また，後述の2次元例は，本論文で構成した具体例として計算する．
% 執筆者確認：実際にセミナーで学んだ範囲と，HHLに関心を持った経緯を補う．

線形方程式
\begin{equation}
  A\bm{x}=\bm{b}
  \label{eq:introduction-linear-system}
\end{equation}
は，線形代数における基本的な問題である．
ここで，$A$は$N$次正方行列，$\bm{b}$は$N$次元ベクトルであり，
$\bm{x}$が求めるべき未知ベクトルである．
行列$A$が可逆であれば，式\eqref{eq:introduction-linear-system}の解は
$\bm{x}=A^{-1}\bm{b}$によって一意に定まる．
線形方程式はそれ自体が重要であるだけでなく，最小二乗法をはじめとする
さまざまな計算を記述する際にも現れる．

Harrow，Hassidim，Lloydが提案したHHLアルゴリズムは，
一定の仮定の下で線形方程式を量子計算によって扱うアルゴリズムである
\cite{shimada2020}．
ただし，HHLアルゴリズムが出力するものは，
$A^{-1}\bm{b}$の成分をすべて並べた古典的なベクトルではない．
非零ベクトル$\bm{b}$の長さを1にそろえ，量子状態$\ket{b}$として与えたとき，
HHLアルゴリズムの目標は，理想的には
\begin{equation}
  \ket{x}
  =
  \frac{A^{-1}\ket{b}}
       {\lVert A^{-1}\ket{b}\rVert}
  \label{eq:introduction-solution-state}
\end{equation}
という量子状態を生成することである．
したがって，HHLアルゴリズムを理解するためには，
線形方程式の解だけでなく，量子状態の変換と測定を考える必要がある．

HHLアルゴリズムは，単一の量子操作によって
式\eqref{eq:introduction-solution-state}を直接生成するものではない．
まず，エルミート行列$A$の固有値を量子位相推定によって量子レジスタに記録する．
次に，固有値に応じた制御回転を補助量子ビットへ作用させ，
固有値$\lambda$の逆数$1/\lambda$を補助量子ビットの振幅に取り入れる．
最後に，逆量子位相推定によって途中計算に用いたレジスタを初期状態へ戻し，
補助量子ビットを測定する．
所定の測定結果が得られた場合に，残された量子状態が
式\eqref{eq:introduction-solution-state}となる．
これらの操作は，量子フーリエ変換，量子位相推定，制御回転，
行列の指数関数，振幅増幅といった複数の数学的事項に基づいている
\cite{shimada2020}．

本論文の目的は，HHLアルゴリズムを構成するこれらの事項を順に定義し，
その性質を証明する．その上で，各段階の量子状態を数式で順に計算する．
本論文では，固有値の逆数が量子状態の振幅に取り入れられる仕組みに注目する．
この仕組みを追うことで，個別に導入される量子計算の基礎事項と，
一つのアルゴリズムの正しさとの関係を具体的に説明できる．
そのため，対象をHHLとその理解に必要な準備に絞る．
特に，次の三点を明らかにする．
\begin{enumerate}[(1)]
  \item 量子状態，測定，複合系および量子ゲートを，
        複素数を成分にもつベクトルと行列を用いて表す．
  \item 量子位相推定など，HHLで使う各処理について，
        必要な仮定を書き，状態がどのように変化するかを計算する．
  \item HHLアルゴリズムの各段階を計算し，
        成功時の出力が式\eqref{eq:introduction-solution-state}となることを証明する．
\end{enumerate}

また，固有値が有限の量子ビット数で正確に表される理想的な場合と，
位相推定や制御回転に近似を含む場合とを区別する．
理想的な場合には出力状態の正しさを詳しく証明し，
近似を含む場合には固有値の誤差が逆数へどのように伝わるかを説明する．
さらに，成功確率，行列の条件数，入力状態の準備，
行列の指数関数に対応する量子操作，出力の読み出し方を整理する．
これにより，HHLの計算手順と，計算量を考えるために必要な条件を
分けて理解することを目指す．
